\documentclass{article}
\usepackage{graphicx} % Required for inserting images
\usepackage[a4paper, total={6.5in, 9in}]{geometry}

\newcommand{\horrule}[1]{\rule{\linewidth}{#1}} %
\usepackage{amsmath,vhistory}

\usepackage[colorlinks=true,linkcolor=blue]{hyperref}%

\usepackage{xcolor,makecell}
\usepackage{hyperref}
\date{\today}

\begin{document}

\date{\today}
\begin{titlepage}
   \vspace*{\stretch{1.0}}
   \begin{center}
        \horrule{2pt} \\[0.5cm] % Thick bottom horizontal rule
              \Large\textbf{PyActiveStorage}\\
              \Large\text{Preliminary Testing Results}\\
        \horrule{2pt} \\[0.5cm] % Thick bottom horizontal rule
         
         \large\text{Varsiha Sothilingam}\\
         \large\text{NCAS-CMS, University of Reading}\\
         \large\text{varsiha.sothilingam@ncas.ac.uk}\\
         \large\text{\today}\\
         \large\text{Version \vhCurrentVersion : Distributed via Slack}\\
   \end{center}
   \vspace*{\stretch{2.0}}





\end{titlepage}
   
\begin{versionhistory}
\vhEntry{1}{25.03.2026}{V.S}{Initial document created with draft of an overview in Section 1. 
                Tests completed so far filled into section 2. An attempt to interpt the results is undergone in Section 3, currently on hold until the last tests are completed. Distributed via Slack}
\vhEntry{2}{\today}{V.S}{hopefully some more tests and the comparisons}
\end{versionhistory}

\setcounter{table}{0}

\section{Overview}

Testing undertaken to demonstrate performance of \texttt{PyActiveStorage}, and to quantify factors which impact its ability to perform remote reductions. Different settings are compared and tested, listed below:
\begin{itemize}
    \item Chunk size
    \item Network: UoR Eduroam vs UoR Ethernet
    \item Methods (min max count sum mean) \textcolor{red}{[Not completed, Note: weights for sum and mean]}
    \item Data retrieval: Object \texttt{s3} store vs \texttt{https} protocol
    \item Maximum threads set \textcolor{red}{[Not completed]}
    \item Cache clearing \textcolor{red}{[Not completed]}
    \end{itemize}

\noindent The performance is evaluated using the following parameters of interest:

\begin{itemize}
    \item Real time: wall-clock
    \item User time: CPU time spent executing process
    \item Sys time: CPU time spent in kernel/OS
    \item Idle time: $t_\text{wait} = t_\text{real} - (t_\text{user} + t_\text{sys})$, where reductionist to receive, respond and talk to disk  
    \item CPU of process: $\frac{t_\text{user}+ t_\text{sys}}{t_\text{real}} \times 100~[\%]$
    \item Memory usage: from max. resident set size: $\frac{\text{RES}}{1024^2}~[\text{MB}]$
\end{itemize}

Scripts and data can be found in the following \href{https://github.com/valeriupredoi/testing_PyActiveStorage/tree/main/tests_Varsiha}{Git repository}.

\newpage



\section{Test Conditions \& Results}

In order to understand the performance, multiple iterations of a given test are completed in order to fill histograms and evaluate the average performance. Caching of the data will directly influence its performance. Therefore, an initial ``cold" run is ran such that for the following ``warm" runs, the caching of the data is already present on the Reductionist. Typically 60 warm runs are recorded, however if a process was very slow, it was capped at an earlier number of iterations. This will be indicated in the respective testing sections. Unless stated otherwise, tests are performed with the number of maximum threads is set to 100, the method applied is \texttt{min} and the data file is retrieved from the \texttt{s3} object store.


%-----------------------------------------------------------------------------------
%                          TEST 1
%-----------------------------------------------------------------------------------
\subsection{Test 1}
\begin{table}[h!] 
    \centering 
    \begin{tabular}{lc}  \hline
                                        & Test 1                                          \\ \hline
        \textbf{File}                   & \texttt{ch330a.pc19790301-bnl.nc}  (17.96  GB)  \\  
        \textbf{Dataset}                & \texttt{UM\_m01s16i202\_vn1106}                 \\ 
        \textbf{Shape}                  & (40, 1920, 2560)                                \\ 
        \textbf{Chunk shape}            & (4, 113, 128)                                   \\ 
        \textbf{Chunks per axis}        & [10, 17, 20]                                    \\ 
        \textbf{Total number of chunks} & 3400                                            \\ 
        \textbf{Slicing applied}        & [0:4, 0:30, 0:30]                               \\ 
        \textbf{No. of Iterations}      & 60                                              \\  \hline
    \end{tabular} 
    \caption{} 
    \label{tab:T1_Conditions} 
\end{table}


\begin{table}[h!]
    \centering
    \begin{tabular}{cc}
        \hline
        \textbf{Results} & Test 1              \\ \hline
        Real Time [s]    & 14.696 $\pm$ 1.842   \\
        User Time [s]    & 1.610  $\pm$ 0.445   \\
        Sys Time [s]     & 0.902  $\pm$ 0.232   \\
        Wait Time [s]    & 12.185 $\pm$ 1.189   \\
        Memory [MB]      & 286.142 $\pm$ 10.856 \\
        CPU Util [\%]    & 16.904 $\pm$ 1.700   \\
        \hline
    \end{tabular}
    \caption{Test 1 Results}
    \label{tab:T1_Results}
\end{table}

%-----------------------------------------------------------------------------------
%                          TEST 2
%-----------------------------------------------------------------------------------
\subsection{Test 2}
\begin{table}[h!] 
    \centering 
    \begin{tabular}{lc}  \hline
                                        & Test 2                                          \\ \hline
        \textbf{File}                   & \texttt{ch330a.pc19790301-bnl.nc}  (17.96  GB)  \\  
        \textbf{Dataset}                &  \texttt{UM\_m01s30i202\_vn1106}                \\ 
        \textbf{Shape}                  &  (40, 7, 1921, 2560)                            \\ 
        \textbf{Chunk shape}            &  (4, 1, 113, 128)                               \\ 
        \textbf{Chunks per axis}        &  [10, 7, 17, 20]                                \\ 
        \textbf{Total number of chunks} &    23800                                        \\ 
        \textbf{Slicing applied}        & [0:4, 0:1, 0:30, 0:30]                          \\ 
        \textbf{No. of Iterations}      &  20                                             \\  \hline
    \end{tabular} 
    \caption{} 
    \label{tab:T2_Conditions} 
\end{table}


\begin{table}[h!]
    \centering
    \begin{tabular}{cc}
        \hline
        \textbf{Results} & Test 2              \\ \hline
        Real Time [s]    & 65.147 $\pm$ 1.622  \\
        User Time [s]    & 6.154  $\pm$ 0.184  \\
        Sys Time [s]     & 3.503  $\pm$ 0.095  \\
        Wait Time [s]    & 55.489 $\pm$ 1.380  \\
        Memory [MB]      & 304.298 $\pm$ 8.933 \\
        CPU Util [\%]    & 14.824 $\pm$ 0.178  \\
        \hline
    \end{tabular}
    \caption{Test 2 Results}
    \label{tab:T2_Results}
\end{table}

%-----------------------------------------------------------------------------------
%                          TEST 3
%-----------------------------------------------------------------------------------
\break
\subsection{Test 3}

\begin{table}[h!] 
    \centering 
    \begin{tabular}{lc}  \hline
                                        & Test 3                                          \\ \hline
        \textbf{File}                   & \texttt{ch330a.pc19790301-def.nc} (17.81  GB)   \\  
        \textbf{Dataset}                &  \texttt{UM\_m01s16i202\_vn1106}  \\ 
        \textbf{Shape}                  &  (40, 1920, 2560)                 \\ 
        \textbf{Chunk shape}            &  (10, 480, 640)                   \\ 
        \textbf{Chunks per axis}        &  [4, 4, 4]                        \\ 
        \textbf{Total number of chunks} &  64                               \\ 
        \textbf{Slicing applied}        &  [0:4, 0:30, 0:30]                \\ 
        \textbf{No. of Iterations}      &  60                               \\  \hline
    \end{tabular} 
    \caption{} 
    \label{tab:T3_Conditions} 
\end{table}


\begin{table}[h!]
    \centering
    \begin{tabular}{cc}
        \hline
        \textbf{Results} & Test 3              \\ \hline
        Real Time [s]    &  4.937 $\pm$ 1.283   \\
        User Time [s]    &  0.657 $\pm$ 0.219   \\
        Sys Time [s]     &  0.345 $\pm$ 0.107   \\
        Wait Time [s]    &  3.935 $\pm$ 0.985   \\
        Memory [MB]      &  269.209 $\pm$ 9.498 \\
        CPU Util [\%]    &  20.136 $\pm$ 1.744  \\
        \hline
    \end{tabular}
    \caption{Test 3 Results}
    \label{tab:T3_Results}
\end{table}

%-----------------------------------------------------------------------------------
%                          TEST 4
%-----------------------------------------------------------------------------------
\subsection{Test 4}
\begin{table}[h!] 
    \centering 
    \begin{tabular}{lc}  \hline
                                        & Test 4                                          \\ \hline
        \textbf{File}                   & \texttt{ch330a.pc19790301-def.nc} (17.81  GB)  \\  
        \textbf{Dataset}                & \texttt{UM\_m01s30i202\_vn1106}     \\ 
        \textbf{Shape}                  & (40, 7, 1921, 2560)                 \\ 
        \textbf{Chunk shape}            & (8, 1, 385, 512)                    \\ 
        \textbf{Chunks per axis}        & [5, 7, 5, 5]                        \\ 
        \textbf{Total number of chunks} & 875                                 \\ 
        \textbf{Slicing applied}        & 0:4, 0:1, 0:30, 0:30]               \\ 
        \textbf{No. of Iterations}      & 20                                  \\  \hline
    \end{tabular} 
    \caption{} 
    \label{tab:T4_Conditions} 
\end{table}


\begin{table}[h!]
    \centering
    \begin{tabular}{cc}
        \hline
        \textbf{Results} & Test 4               \\ \hline
        Real Time [s]    &  18.605 $\pm$ 0.816  \\
        User Time [s]    &  1.876 $\pm$ 0.135   \\
        Sys Time [s]     &  1.058 $\pm$ 0.054   \\
        Wait Time [s]    &  15.673 $\pm$ 0.654  \\
        Memory [MB]      &  286.00 $\pm$ 9.291  \\
        CPU Util [\%]    &  15.756 $\pm$ 0.463  \\
        \hline
    \end{tabular}
    \caption{Test 4 Results}
    \label{tab:T4_Results}
\end{table}

%-----------------------------------------------------------------------------------
%                          TEST 5
%-----------------------------------------------------------------------------------
\break
\subsection{Test 5}
\begin{table}[h!] 
    \centering 
    \begin{tabular}{lc}  \hline
                                        & Test 5                                                                                                        \\ \hline
        \textbf{File}                   &  \makecell{\texttt{cl\_Amon\_UKESM1-0-LL\_ssp370SST}\\ \texttt{-lowNTCF\_r1i1p1f2\_gn\_205001-209912.nc} (2.36 GB)}     \\  
        \textbf{Dataset}                &\texttt{cl}                                                                                \\ 
        \textbf{Shape}                  &  (600, 85, 144, 192)                                                                                              \\ 
        \textbf{Chunk shape}            & (1, 43, 72, 96)                                                                                                 \\ 
        \textbf{Chunks per axis}        & (600, 2, 2, 2)                                                                                                  \\ 
        \textbf{Total number of chunks} & 4800                                                                                                          \\ 
        \textbf{Slicing applied}        & [0:4, 0:1, 0:30, 0:30]                                                                                             \\ 
        \textbf{No. of Iterations}      & 20                                                                                                            \\  \hline
    \end{tabular} 
    \caption{} 
    \label{tab:T5_Conditions} 
\end{table}


\begin{table}[h!]
    \centering
    \begin{tabular}{cc}
        \hline
        \textbf{Results} & Test 5              \\ \hline
        Real Time [s]    & 30.014 $\pm$ 0.898  \\
        User Time [s]    & 2.981 $\pm$ 0.098   \\
        Sys Time [s]     & 1.606 $\pm$ 0.073   \\
        Wait Time [s]    & 25.428 $\pm$ 0.772  \\
        Memory [MB]      & 289.440 $\pm$ 8.367 \\
        CPU Util [\%]    & 15.282 $\pm$ 0.266  \\
        \hline
    \end{tabular}
    \caption{Test 5 Results}
    \label{tab:T5_Results}
\end{table}

%-----------------------------------------------------------------------------------
%                          TEST 6
%-----------------------------------------------------------------------------------
\subsection{Test 6}

\begin{table}[h!] 
    \centering 
    \begin{tabular}{lc}  \hline
                                        & Test 6                                         \\ \hline
        \textbf{File}                   & \texttt{da193a\_25\_day\_\_198808-198808.nc}   \\  
        \textbf{Dataset}                &\texttt{m01s06i247\_4 }                         \\ 
        \textbf{Shape}                  & (30, 39, 325, 432)                             \\ 
        \textbf{Chunk shape}            & (1, 39, 325, 432)                              \\ 
        \textbf{Chunks per axis}        &  (30, 1, 1, 1) )                               \\ 
        \textbf{Total number of chunks} & 30                                             \\ 
        \textbf{Slicing applied}        & [0:4, 0:1, 0:30, 0:30]                         \\ 
        \textbf{No. of Iterations}      & 60                                             \\  \hline
    \end{tabular}
    \caption{} 
    \label{tab:T6_Conditions} 
\end{table}

\begin{table}[h!]
    \centering
    \begin{tabular}{cc}
        \hline
        \textbf{Results} & Test 6              \\ \hline
        Real Time [s]    & 1.242 $\pm$ 0.091  \\
        User Time [s]    & 0.333 $\pm$ 0.012  \\
        Sys Time [s]     & 0.138 $\pm$ 0.008  \\
        Wait Time [s]    & 0.772 $\pm$ 0.087  \\
        Memory [MB]      & 206.392 $\pm$ 1.525\\
        CPU Util [\%]    & 38.032 $\pm$ 2.294 \\
        \hline
    \end{tabular}
    \caption{Test 6 Results}
    \label{tab:T6_Results}
\end{table}


%-----------------------------------------------------------------------------------
%                          TEST 7
%-----------------------------------------------------------------------------------
\break
\subsection{Test 7}

\begin{table}[h!] 
    \centering 
    \begin{tabular}{lc}  \hline
                                        & Test 7                                         \\ \hline
        \textbf{File}                   & \texttt{da193a\_25\_6hr\_t\_pt\_cordex\_\_198807-198807.nc} (6.13 GB)   \\  
        \textbf{Dataset}                &\texttt{ m01s30i111  }                         \\ 
        \textbf{Shape}                  & (120, 85, 324, 432)                             \\ 
        \textbf{Chunk shape}            & (1, 85, 324, 432)                              \\ 
        \textbf{Chunks per axis}        &  (120, 1, 1, 1) )                               \\ 
        \textbf{Total number of chunks} & 120                                             \\ 
        \textbf{Slicing applied}        & [0:4, 0:1, 0:30, 0:30]                         \\ 
        \textbf{No. of Iterations}      & 20                                             \\  \hline
    \end{tabular}
    \caption{} 
    \label{tab:T7_Conditions} 
\end{table}

\begin{table}[h!]
    \centering
    \begin{tabular}{cc}
        \hline
        \textbf{Results} & Test 7              \\ \hline
        Real Time [s]    & 2.917 $\pm$ 0.167  \\
        User Time [s]    & 0.497 $\pm$ 0.052  \\
        Sys Time [s]     & 0.232 $\pm$ 0.031  \\
        Wait Time [s]    & 2.187 $\pm$ 0.106  \\
        Memory [MB]      & 276.494 $\pm$ 8.837\\
        CPU Util [\%]    & 24.957 $\pm$ 1.729 \\
        \hline
    \end{tabular}
    \caption{Test 7 Results}
    \label{tab:T7_Results}
\end{table}

%-----------------------------------------------------------------------------------
%                          TEST 8
%-----------------------------------------------------------------------------------
\break
\subsection{Test 8}
Test 8 has the same conditions as Test 1 except now it is ran on the UoR Eduroam network instead in order to compare the difference between the network and wifi performance of the PyActiveStorage. As running on the WiFi is much slower, only 10 iterations are ran to determine the results.

\begin{table}[h!]
    \centering
    \begin{tabular}{cc}
        \hline
        \textbf{Results} & Test 8              \\ \hline
        Real Time [s]    &80.063 $\pm$ 18.940 \\
        User Time [s]    &6.533 $\pm$ 0.672   \\
        Sys Time [s]     &3.132 $\pm$ 0.390   \\
        Wait Time [s]    &70.398 $\pm$ 19.189 \\
        Memory [MB]      &290.321 $\pm$ 5.774 \\
        CPU Util [\%]    &12.533 $\pm$ 2.230  \\
        \hline
    \end{tabular}
    \caption{Test 8 Results}
    \label{tab:T8_Results}
\end{table}



%-----------------------------------------------------------------------------------
%                          TEST 9 & 10
%-----------------------------------------------------------------------------------

\subsection{Test 9}
Test 9 is the same as Test 5 except in how the data is read in. Test 5 uses the s3 object store, with the file stored in the bnl bucket. For Test 9, the same file is accessed via https of the ESGF CEDA catalogue . 

\begin{table}[h!]
    \centering
    \begin{tabular}{cc}
        \hline
        \textbf{Results} & Test 9          \\ \hline
        Real Time [s]    &  95.014 $\pm$ 13.666              \\
        User Time [s]    &  3.030 $\pm$ 0.115                \\
        Sys Time [s]     &  1.432 $\pm$ 0.045                \\
        Wait Time [s]    &  90.552 $\pm$ 13.556              \\
        Memory [MB]      &  171.616 $\pm$ 4.515              \\
        CPU Util [\%]    &  4.771 $\pm$ 0.554                \\
        \hline
    \end{tabular}
    \caption{Test 9 Results}
    \label{tab:T9_Results}
\end{table}




\section{Tests \textcolor{red}{Rough draft, need to check still}}


\subsection{File sizes, data shapes and chunking variations}
\subsubsection{Data Shape}
This test aims to understand how the choice of a HDF5 dataset, i.e its shape, influences the performance. Shown in Table \ref{tab:T1T2_Conditions} are the difference in the datasets chosen. The additional dimension in the \texttt{UM\_m01s30i202\_vn1106} dataset slows down the running time and therefore it is tested on a smaller number of iterations. 
\begin{table}[h!] 
    \centering 
    \begin{tabular}{lcc}  \hline
                                        & Test 1                            & Test 2                              \\ \hline
        \textbf{File}                   & \multicolumn{2}{c}{ \texttt{ch330a.pc19790301-bnl.nc}  (17.96  GB)}     \\                                       
        \textbf{Dataset}                &\texttt{UM\_m01s16i202\_vn1106}    & \texttt{UM\_m01s30i202\_vn1106}     \\ 
        \textbf{Shape}                  & (40, 1920, 2560)                  & (40, 7, 1921, 2560)                 \\ 
        \textbf{Chunk shape}            & (4, 113, 128)                     & (4, 1, 113, 128)                    \\ 
        \textbf{Chunks per axis}        & [10, 17, 20]                      & [10, 7, 17, 20]                     \\ 
        \textbf{Total number of chunks} & 3400                              &   23800                             \\ 
        \textbf{Slicing applied}        & [0:4, 0:30, 0:30]                 &[0:4, 0:1, 0:30, 0:30]               \\ 
        \textbf{No. of Iterations}      & 60                                & 20                                  \\  \hline
    \end{tabular} 
    \caption{Properties of the tests performed when varying the HDF5 Dataset in the \texttt{ch330a.pc19790301-bnl.nc} file. For both tests the number of maximum threads is set to 100, the method applied is \texttt{min} and the data file is retrieved from the \texttt{s3} object store.} 
    \label{tab:T1T2_Conditions} 
\end{table}



\begin{table}[h!] 
    \centering 
    \begin{tabular}{lccc}
        \hline
        \textbf{Results}  & Test 1               & Test 2               & $T1/T2$ \\ \hline
        Real Time [s]     & 14.696 $\pm$ 1.842   & 65.147 $\pm$ 1.622   & 0.226 $\pm$ 0.029 \\
        User Time [s]     & 1.610  $\pm$ 0.445   & 6.154  $\pm$ 0.184   & 0.262 $\pm$ 0.073 \\
        Sys Time [s]      & 0.902  $\pm$ 0.232   & 3.503  $\pm$ 0.095   & 0.257 $\pm$ 0.067 \\
        Wait Time [s]     & 12.185 $\pm$ 1.189   & 55.489 $\pm$ 1.380   & 0.220 $\pm$ 0.022 \\
        Memory [MB]       & 286.142 $\pm$ 10.856 & 304.298 $\pm$ 8.933  & 0.940 $\pm$ 0.046 \\
        CPU Util [\%]     & 16.904 $\pm$ 1.700   & 14.824 $\pm$ 0.178   & 1.140 $\pm$ 0.116 \\
        \hline
    \end{tabular}
    \caption{Results obtained from Test 1 and 2, showing all parameters of interest. The ratio between the two tests, $T1/T2$ is shown with propagated uncertainties.} 
    \label{tab:T1T2_Results} 
\end{table}

\noindent Table \ref{tab:T1T2_Results} shows the averaged results for the parameters of interest. Main observations are as follows:
\begin{itemize}
    \item CPU and Memory values are compatible between the tests
    \item Key difference in the Wait Time (I/O of the data). Although only one additional chunk is selected, the variable has seven times more chunks in the file which will increase the overhead due to factors such as additional metadata lookup and chunk indexing. 

\end{itemize}

\subsubsection{Coarser Chunking}
The same HDF5 datasets are studied in another file where the number of chunks is coarser. The number of chunks along with the testing conditions can be found in Table \ref{tab:T3T4_Conditions}. 

\begin{table}[h!] 
    \centering 
    \begin{tabular}{lcc}  \hline
                                        & Test 3                            & Test 4                              \\ \hline
        \textbf{File}                   & \multicolumn{2}{c}{ \texttt{ch330a.pc19790301-def.nc} (17.81  GB)}      \\                                       
        \textbf{Dataset}                & \texttt{UM\_m01s16i202\_vn1106}   & \texttt{UM\_m01s30i202\_vn1106}     \\ 
        \textbf{Shape}                  & (40, 1920, 2560)                  & (40, 7, 1921, 2560)                 \\ 
        \textbf{Chunk shape}            & (10, 480, 640)                    & (8, 1, 385, 512)                    \\ 
        \textbf{Chunks per axis}        & [4, 4, 4]                         & [5, 7, 5, 5]                        \\ 
        \textbf{Total number of chunks} & 64                                & 875                                 \\ 
        \textbf{Slicing applied}        & [0:4, 0:30, 0:30]                 &[0:4, 0:1, 0:30, 0:30]               \\ 
        \textbf{No. of Iterations}      & 60                                & 20                                  \\  \hline
    \end{tabular} 
    \caption{Variables compared from the \texttt{ch330a.pc19790301-def.nc} file, for a maximum number of threads of 100, method applied is \texttt{min} and the data file is retrieved from the \texttt{s3} object store.} 
    \label{tab:T3T4_Conditions} 
\end{table}

\begin{table}[h!]
    \centering
    \begin{tabular}{cccc}
        \hline
        \textbf{Results} & Test 3               & Test 4               & $T3/T4$ \\ \hline
        Real Time [s]    & 4.937 $\pm$ 1.283    & 18.605 $\pm$ 0.816   & 0.265 $\pm$ 0.070 \\
        User Time [s]    & 0.657 $\pm$ 0.219    & 1.876 $\pm$ 0.135    & 0.350 $\pm$ 0.118 \\
        Sys Time [s]     & 0.345 $\pm$ 0.107    & 1.058 $\pm$ 0.054    & 0.326 $\pm$ 0.103 \\
        Wait Time [s]    & 3.935 $\pm$ 0.985    & 15.673 $\pm$ 0.654   & 0.251 $\pm$ 0.064 \\
        Memory [MB]      & 269.209 $\pm$ 9.498  & 286.00 $\pm$ 9.291   & 0.941 $\pm$ 0.047 \\
        CPU Util [\%]    & 20.136 $\pm$ 1.744   & 15.756 $\pm$ 0.463  & 1.278 $\pm$ 0.115 \\
        \hline
    \end{tabular}
    \caption{Variables compared from the \texttt{ch330a.pc19790301-def.nc} file. The factor $T3/T4$ is shown with propagated uncertainty based on the relative standard deviations of both tests.}
    \label{tab:T3T4_Results}
\end{table}

\break
\noindent Table \ref{tab:T3T4_Results} shows the averaged results for the parameters of interest. Main observations are as follows:
\begin{itemize}
    \item The ratios for all parameters are very similar to that observed in Table \ref{tab:T1T2_Results}, showing that the additional dimension in the data shape has the same impact, even if other axis are chunked coarser. 
\end{itemize}

\begin{table}[h!] 
    \centering 
    \begin{tabular}{lcc}  \hline
        \textbf{Results} & Test 1               &   Test 3              \\ \hline
        Real Time [s]    & 14.696 $\pm$ 1.842   &   4.937 $\pm$ 1.283   \\
        User Time [s]    & 1.610  $\pm$ 0.445   &   0.657 $\pm$ 0.219   \\
        Sys Time [s]     & 0.902  $\pm$ 0.232   &   0.345 $\pm$ 0.107   \\
        Wait Time [s]    & 12.185 $\pm$ 1.189   &   3.935 $\pm$ 0.985   \\
        Memory [MB]      & 286.142 $\pm$ 10.856 &   269.209 $\pm$ 9.498 \\
        CPU Util [\%]    & 16.904 $\pm$ 1.700   &   20.136 $\pm$ 1.744  \\ \hline \hline


        \textbf{Results}   & Test 2               & Test 4             \\ \hline
        Real Time [s]      & 65.147 $\pm$ 1.622   & 18.605 $\pm$ 0.816 \\
        User Time [s]      & 6.154  $\pm$ 0.184   & 1.876 $\pm$ 0.135  \\
        Sys Time [s]       & 3.503  $\pm$ 0.095   & 1.058 $\pm$ 0.054  \\
        Wait Time [s]      & 55.489 $\pm$ 1.380   & 15.673 $\pm$ 0.654 \\
        Memory [MB]        &304.298 $\pm$ 8.933   & 286.00 $\pm$ 9.291 \\
        CPU Util [\%]      & 14.824 $\pm$ 0.178   & 15.756 $\pm$ 0.463 \\ \hline \hline

        
    \end{tabular} 
    \caption{Variables compared from the \texttt{ch330a.pc19790301-def.nc} file, for a maximum number of threads of 100, method applied is \texttt{min} and the data file is retrieved from the \texttt{s3} object store.} 
    \label{tab:T1T2T3T4_ChunkResults}
\end{table}


\noindent The impact of the coarser chunking can be compared for the same dataset between the two files. The results for how the chunking impacts the parameters of interest are shown in Table \ref{tab:T1T2T3T4_ChunkResults} where Test 1 and 3 compare for the \texttt{UM\_m01s16i202\_vn1106} data set, and Test 2 and 4 the \texttt{UM\_m01s30i202\_vn1106} data set. The main observations are as follows:

\begin{itemize}
    \item The smaller number of chunks allows for the reductions to be completed faster.
    \item The CPU usuage is similar due to a short process time, there are fewer idle resources.
    \item T2/T4 ratio takes more time due to 
\end{itemize}


\subsection{UoR Ethernet vs Eduroam}






\subsection{s3 v https}








\end{document}
